%!TEX root =  tfg.tex
\chapter{Contexto}

\begin{quotation}[Philosopher]{Plato (c. 427--347 BC)}
The beginning is the most important part of the work.
\end{quotation}

\begin{abstract}
En este capítulo se describe el marco contextual de \textit{Via Inferni}. Se presenta la
relevancia del videojuego independiente, se resumen las características principales del
género \textit{roguelike}, se revisan títulos de referencia y se explica el papel de la
\textit{Divina Comedia} como base temática del proyecto.
\end{abstract}

\section{El mundo de los videojuegos}

La industria del videojuego se ha consolidado como uno de los sectores de entretenimiento
más relevantes de las últimas décadas. Su crecimiento no se limita a las grandes
producciones, sino que también incluye proyectos independientes capaces de alcanzar una
visibilidad significativa mediante propuestas más acotadas, originales y ajustadas a
equipos reducidos.

Este contexto resulta especialmente importante para un TFG como \textit{Via Inferni}. El
proyecto no pretende competir con una producción profesional de gran escala, sino adoptar
un enfoque cercano al desarrollo independiente: alcance controlado, mecánicas bien
definidas y aprovechamiento de herramientas accesibles. En este sentido, motores como
Unity han facilitado que equipos pequeños puedan crear videojuegos completos con una base
técnica sólida y una curva de aprendizaje asumible \cite{unityLearnGettingStarted}.

\section{El género roguelike}

El género \textit{roguelike} tiene su origen en los juegos de rol por texto de los años
ochenta, aunque con el tiempo se ha adaptado a múltiples estilos, plataformas y ritmos de
juego. Su interés para este proyecto reside en que permite generar experiencias variables
sin diseñar manualmente cada partida completa, lo que encaja con las limitaciones de
tiempo y recursos de un trabajo académico.

En términos generales, un \textit{roguelike} suele apoyarse en tres elementos:

\begin{itemize}
    \item \textbf{Generación procedural de contenido:} los niveles se construyen mediante
    algoritmos, de forma que cada partida presenta una disposición distinta de salas,
    enemigos, recompensas y riesgos.
    
    \item \textbf{Reinicio tras la derrota:} la muerte del personaje devuelve al jugador
    al inicio del recorrido o a un estado muy anterior, lo que da importancia a cada
    decisión y favorece la repetición de partidas.
    
    \item \textbf{Aprendizaje progresivo del jugador:} aunque el avance de una partida
    pueda perderse, el jugador mejora al conocer patrones, enemigos, objetos y estrategias.
\end{itemize}

En \textit{Via Inferni}, estos elementos se adaptan a una experiencia 2D de acción. La
generación de salas sostiene la variabilidad, el combate en tiempo real aporta tensión y
el avance por círculos permite introducir una dificultad creciente sin depender
únicamente de la aleatoriedad.

\section{Estado del arte}

A lo largo de los últimos años, distintos títulos han ampliado el significado práctico del
término \textit{roguelike}. Algunos mantienen la tradición clásica, mientras que otros
incorporan acción en tiempo real, narrativa persistente o sistemas de progresión más
accesibles. Para este proyecto se han tomado como referencia los siguientes casos:

\begin{itemize}
    \item \textbf{Rogue} (1980): título que dio nombre al género y estableció elementos
    básicos como la exploración de mazmorras, la generación procedural y la muerte
    permanente.
    
    \item \textbf{NetHack} (1987-actualidad): evolución de los roguelikes clásicos, útil
    como referencia por la profundidad de sus sistemas y por la gran cantidad de
    situaciones emergentes que permite.
    
    \item \textbf{Dungeon Crawl: Stone Soup} (2006-actualidad): ejemplo de diseño centrado
    en estrategia, exploración y toma de decisiones, manteniendo una estructura clásica de
    mazmorra.
    
    \item \textbf{The Binding of Isaac} (2011): roguelike de acción en perspectiva cenital
    \cite{bindingIsaacOfficial}. Resulta especialmente relevante por su estructura de
    salas, su sistema de objetos y su forma de combinar dificultad, azar y partidas
    cortas.
    
    \item \textbf{Hades} (2020): roguelike de acción que integra progresión narrativa con
    repetición de intentos \cite{hadesOfficial}. Su interés para este trabajo está en la
    relación entre estructura cíclica, personajes y avance argumental.
\end{itemize}

Estas referencias muestran que el género puede abordarse desde distintos niveles de
complejidad. En el caso de \textit{Via Inferni}, se toma una aproximación acotada:
generación de salas conectadas, progresión por zonas, enemigos diferenciados y sistemas de
combate y recompensa suficientes para sostener la rejugabilidad del prototipo.

\section{El Infierno de Dante como base narrativa}

Mientras que el género \textit{roguelike} proporciona la estructura jugable, la
\textit{Divina Comedia} aporta el marco simbólico y temático del proyecto. El Infierno de
Dante se organiza en círculos asociados a distintos pecados, castigos y figuras, lo que
permite convertir una obra literaria clásica en una estructura comprensible para el diseño
de niveles.

Esta adaptación no busca reproducir la obra de forma literal. Su objetivo es utilizarla
como inspiración para definir ambientaciones, enemigos, progresión y tono general. De este
modo, cada círculo puede funcionar como una zona con identidad propia, manteniendo al mismo
tiempo una coherencia global dentro del recorrido del jugador.

\begin{itemize}
    \item \textbf{Estructura de niveles:} los círculos permiten estructurar el juego en
    zonas progresivas con dificultad y ambientación diferenciadas.
    
    \item \textbf{Identidad visual y temática:} cada pecado ofrece una base para diseñar
    enemigos, escenarios y recompensas de forma coherente.
    
    \item \textbf{Dualidad de personajes:} la relación entre Dante y Virgilio se traslada a
    una mecánica de alternancia entre personajes con capacidades complementarias.
    
    \item \textbf{Diferenciación del prototipo:} el uso de una referencia literaria aporta
    una identidad concreta frente a otros roguelikes centrados únicamente en la mecánica.
\end{itemize}

El valor de esta base narrativa está en que sirve como guía de diseño. Ayuda a delimitar
decisiones artísticas y jugables, y evita que la generación procedural produzca una
experiencia desconectada de la temática general del proyecto.
